% ==================== 2.6. COMPARATIVA DE RESULTADOS ====================
\section{Comparativa de resultados, rendimiento y limitaciones}\label{sec:comparativa}

Para la constitución de este Estado del Arte se ha realizado una revisión bibliográfica sistemática de la cual se han seleccionado doce estudios de alto impacto. A partir de ellos, se ha elaborado una síntesis comparativa de resultados, considerando la patología objetivo, la arquitectura de modelado, las características vocales empleadas y las métricas de rendimiento reportadas. La Tabla~\ref{tab:estudios_neurovoz} resume la información más relevante y permite observar tendencias metodológicas comunes entre los distintos grupos de investigación.

\begin{landscape}
    \begin{table*}[ht]
        \centering
        \footnotesize % Subimos el tamaño, ahora tenemos espacio de sobra
        \renewcommand{\arraystretch}{1.3} % Da un poco de "aire" vertical entre las filas
        \begin{tabularx}{\linewidth}{
            c 
            >{\raggedright\arraybackslash}p{3.5cm} 
            l 
            c 
            >{\raggedright\arraybackslash}X 
            >{\raggedright\arraybackslash}X 
            >{\raggedright\arraybackslash}p{2.6cm} 
            >{\raggedright\arraybackslash}X
        }
            \toprule
            \textbf{Nº} & \textbf{Referencia} & \textbf{Enfermedad} & \textbf{Dataset (N)} & \textbf{Modelo(s)} & \textbf{Características de Entrada} & \textbf{Métricas Clave} & \textbf{Observaciones} \\
            \midrule
            1  & García-Gutiérrez et al. \cite{garcia-gutierrezUnveilingSoundCognitive2024}    & Alzheimer  & 1500    & RF, XGB, SVM, KNN                & Paralingüísticas (voz espontánea) & F1: 84--92\%                  & Enfoque ML clásico; dataset de gran volumen      \\
            2  & Worasawate et al. \cite{worasawate2021}          & Parkinson  & 4051    & CNN (LeNet, ResNet, VGG)         & Espectrogramas /a/ sostenida      & Acc: 97.7--99.3\%             & Grabaciones móviles; alta capacidad de generalización    \\
            3  & Tougui et al. \cite{tougui2020}              & Parkinson  & 1490    & SVM, RF, XGB                     & 138 acústicas y cepstrales        & Acc: 95.8\%; Sens: 95.3\%     & Alto rendimiento con descriptores clásicos  \\
            4  & Moro-Velázquez et al. \cite{moro2020, moro2021new} & Parkinson  & 79--89  & GMM-UBM, PLDA                    & x-/i-vectors, MFCC                & Acc: 85--94\%; AUC: 0.91--0.99 & Embeddings robustos; validación cruzada estricta \\
            5  & Agbavor \& Liang \cite{agbavor2022}           & Alzheimer  & 237     & SVM, RF, LogReg                  & data2vec embeddings               & AUC: 0.835--0.846            & Uso de aprendizaje auto-supervisado                \\
            6  & Dao et al. \cite{dao2025}                 & Parkinson  & --      & Wav2Vec 2.0, Whisper             & Foundation embeddings             & AUC: 91.4\%                  & Representaciones profundas translingüísticas              \\
            7  & Muntaha et al. \cite{muntaha2024}             & Parkinson  & --      & LSTM-BiGRU + EBM                 & Multisource vocal                 & Acc: 96.3--97.5\%             & Modelo híbrido con IA Explicable (XAI)                    \\
            8  & Shen et al. \cite{shen2024}                & Parkinson  & 81      & Hybrid CNN-RNN                   & MFCCs, jitter, shimmer            & Acc: 91.1\%; Sens: 92.5\%     & Alta interpretabilidad clínica con XAI              \\
            9  & Gutz et al. \cite{gutz2019}                & ELA        & 123     & SVM                              & ComParE13 (openSMILE)             & AUC: 0.85--0.99              & Detección de sintomatología bulbar temprana              \\
            10 & He et al. \cite{he2023}                  & Alzheimer  & 119     & RF                               & Acústicas + prosódicas + texto    & Acc: >90\%                   & Enfoque multimodal (acústica y PNL)                      \\
            11 & Álvarez Marquina et al. \cite{alvarez2022}    & Parkinson  & 1280    & CNN                              & Formantes (/a/ sostenida)         & Sens: 96\%; Spec: 99\%        & Alta especificidad diagnóstica                      \\
            12 & Tena et al. \cite{tenaALS}                & ELA        & --      & SVM, RF                          & Componentes cuasi-periódicas      & Acc: 78.7--91\%               & Variabilidad condicionada al dataset                \\
            \bottomrule
        \end{tabularx}
        \caption{Síntesis de la revisión bibliográfica sobre la detección de enfermedades neurodegenerativas mediante Inteligencia Artificial y análisis de voz.}
        \label{tab:estudios_neurovoz}
    \end{table*}
\end{landscape}

\newpage

Como se observa en la Tabla~\ref{tab:estudios_neurovoz}, los estudios centrados en la Enfermedad de Parkinson conforman el grueso de la investigación actual y reportan los techos de rendimiento más elevados (97–99\% de precisión), especialmente cuando se aplican arquitecturas profundas sobre espectrogramas o \textit{embeddings}. En la Enfermedad de Alzheimer, las métricas se sitúan en una horquilla del 84\% al 92\%, destacando el uso de características paralingüísticas y modelos auto-supervisados. Por su parte, la varianza en los estudios de Esclerosis Lateral Amiotrófica (78–99\%) subraya la alta dependencia del conjunto de datos. Estas cifras, si bien resultan extraordinariamente prometedoras, exigen una interpretación cautelosa debido a la heterogeneidad metodológica.

\subsection{Resultados desglosados por patología}

\subsubsection*{Enfermedad de Parkinson (EP)}
Constituye la afección neurodegenerativa más modelada computacionalmente. Las precisiones globales oscilan entre el 85\% y el 99\%, reportándose áreas bajo la curva (AUC) superiores a 0,90 en investigaciones como las de Moro-Velázquez et al. \cite{moro2020, moro2021new} y Worasawate et al. \cite{worasawate2021}. Las Redes Neuronales Convolucionales (CNN) y los modelos fundamentados en \textit{embeddings} exhiben una superioridad clara en corpus de gran volumen, mientras que los algoritmos de \textit{Machine Learning} tradicional (SVM, RF) preservan su eficacia en muestras limitadas sometidas a una selección de características exhaustiva.

\subsubsection*{Enfermedad de Alzheimer (EA)}
La literatura sobre Alzheimer reporta resultados clínicamente sólidos, aunque ligeramente más moderados. García-Gutiérrez et al. \cite{garcia-gutierrezUnveilingSoundCognitive2024} consolidaron puntuaciones F1 de entre el 84\% y el 92\% mediante ensambles clásicos, mientras que Agbavor y Liang \cite{agbavor2022} validaron la idoneidad de los \textit{embeddings} auto-supervisados (data2vec) obteniendo un AUC cercano a 0,84. Se evidencia una tendencia de mejora significativa en aquellos modelos que fusionan la dimensión puramente acústica con el análisis lingüístico \cite{he2023}, al capturar simultáneamente el deterioro motor y el déficit semántico-cognitivo.

\subsubsection*{Esclerosis Lateral Amiotrófica (ELA)}
El volumen de publicaciones orientadas a la ELA es menor, pero sus hallazgos son determinantes. Gutz et al. \cite{gutz2019} demostraron la viabilidad de detectar la afectación bulbar preclínica alcanzando un AUC de 0,85–0,99 mediante el conjunto de descriptores ComParE13. Paralelamente, Tena et al. \cite{tenaALS} reportaron precisiones de hasta el 91\% empleando vectores cuasi-periódicos. La principal barrera en este dominio continúa siendo la limitada disponibilidad de bases de datos públicas de pacientes con ELA.

\subsection{Limitaciones transversales de los estudios}
A pesar de los altos índices de precisión reportados, la revisión exhaustiva de la literatura pone de manifiesto una serie de limitaciones estructurales que condicionan la generalización clínica de los modelos actuales:

\begin{itemize}
    \item \textbf{Escasez y heterogeneidad de los datos:} Las muestras oscilan desde unas pocas decenas hasta varios miles de sujetos, presentando asimetrías severas en idioma, rango de edades y calidad de captura. Esta falta de estandarización obstaculiza el análisis comparativo e induce sesgos demográficos.
    
    \item \textbf{Divergencia en tareas vocales e instrumentación:} La amalgama de protocolos fonatorios (vocales sostenidas, lectura dirigida o habla espontánea) altera fundamentalmente los biomarcadores subyacentes. Asimismo, la variabilidad del hardware (desde micrófonos calibrados en cabinas insonorizadas hasta \textit{smartphones} en entornos ruidosos) compromete la fiabilidad del espectro acústico.
    
    \item \textbf{Validación interna y riesgo de sobreajuste (\textit{Overfitting}):} La inmensa mayoría de las investigaciones restringe su evaluación a técnicas de validación cruzada interna (\textit{K-Fold}), omitiendo la validación sobre corpus de prueba externos e independientes. Este déficit procedimental puede inflar artificialmente las métricas, especialmente en arquitecturas de \textit{Deep Learning} entrenadas con corpus reducidos.
    
    \item \textbf{Inconsistencia métrica y estadística:} La carencia de un estándar de reporte (alternando \textit{Accuracy}, F1-Score, Sensibilidad o AUC) dificulta la evaluación cruzada. Además, es infrecuente la publicación de intervalos de confianza o pruebas de significancia que respalden la solidez de los porcentajes obtenidos.
    
    \item \textbf{Opacidad algorítmica (\textit{Black Box}):} La interpretabilidad sigue constituyendo una barrera para la adopción clínica. Solo investigaciones muy recientes \cite{shen2024, muntaha2024} han comenzado a incorporar técnicas de Inteligencia Artificial Explicable (XAI) para auditar qué descriptores acústicos impulsan las decisiones de la red neuronal.
\end{itemize}

\subsection{Conclusiones de la revisión}
El análisis de la literatura consolida un patrón inequívoco: las arquitecturas de aprendizaje profundo y los modelos fundacionales representan el techo de rendimiento actual en la detección de patologías neurodegenerativas mediante la voz, logrando precisiones que rozan el 99\% en Parkinson y áreas bajo la curva superiores a 0,90 en Alzheimer y ELA. 

Sin embargo, la madurez de la disciplina está supeditada a la resolución de sus vulnerabilidades estructurales. La evidencia actual ratifica el inmenso potencial de la señal de voz como biomarcador digital; no obstante, el avance hacia el diagnóstico clínico real exigirá la adopción de protocolos de captura estandarizados, la evaluación sistemática sobre conjuntos de datos externos y la integración obligatoria de mecanismos de explicabilidad (XAI).